% Figure: steady temperature along a cooling fin for three values of
% the dimensionless fin parameter mL. The profile is the hyperbolic
% cosine of the insulated-tip fin; larger mL means a colder, shorter
% effective fin. Shows how the fin length past which extra material
% buys almost no extra heat rejection is set by mL of order two to
% three.
\begin{figure}[htbp]
\centering
\begin{tikzpicture}
\begin{axis}[
  width=12cm, height=6.5cm,
  xlabel={fractional distance along fin $x/L$},
  ylabel={excess temperature $\theta/\theta_b = (T-T_\infty)/(T_b-T_\infty)$},
  xmin=0, xmax=1,
  ymin=0, ymax=1.05,
  grid=both,
  major grid style={engblack!20},
  minor grid style={engblack!10},
  legend pos=south west,
  legend cell align=left,
  font=\small,
  samples=120,
  domain=0:1,
]
% Insulated-tip fin: theta/theta_b = cosh(mL(1-x/L)) / cosh(mL).
\addplot[engblue, very thick] { cosh(0.5*(1-x)) / cosh(0.5) };
\addlegendentry{$mL = 0.5$ (short/thick fin)}

\addplot[engorange, very thick, dashed] { cosh(2*(1-x)) / cosh(2) };
\addlegendentry{$mL = 2$ (working design)}

\addplot[engred, very thick, dotted] { cosh(5*(1-x)) / cosh(5) };
\addlegendentry{$mL = 5$ (long/thin fin)}

\end{axis}
\end{tikzpicture}
\caption[Steady temperature along a cooling fin]{Steady excess
temperature along an insulated-tip fin,
$\theta/\theta_b = \cosh(mL(1-x/L))/\cosh(mL)$, for three values of
the dimensionless fin parameter $mL = L\sqrt{hP/(kA_c)}$. At $mL =
0.5$ the fin stays near the base temperature over its whole length,
so most of its surface rejects heat effectively. At $mL = 5$ the
temperature has collapsed to the ambient within the first fifth of
the length, so the outer four-fifths carry almost no temperature
excess and reject almost no heat: adding length past $mL \approx 2$
to $3$ buys diminishing return. The knee near $mL \approx 2$ is the
working design target, the same $\tanh(mL)/mL$ efficiency rolloff the
heat-transfer engineer reads from a fin chart.}
\label{fig:vol02:ch12:fin-profile}
\end{figure}
